\SourceLeaf{44}vitur, terminum generalem $r^{\text{tum}} = \cos\left(2r+\frac{1-(-1)^n}{2}\right)\varphi$.
Manifesto enim existente $n$ pari contrahitur haec forma in $\cos 2r\varphi$,
existente impari abit in $\cos(2r+1)\varphi$.

In dissertatiuncula ipsa generalem hunc notandi modum nonnisi in formula
finitima (\S\,11) adhibui. Licet enim iam in disquisitionibus quibus ad
istam pervenitur cum fructu adplicari possit, abstinendum ipso censui, ne
propter addenda quaedam de transmutationibus eiusmodi signorum, si in
calculos ipsos ingrediantur, ad nimiam prolixitatem abreptus,
dissertatiunculae limites excederem. Satius mihi visum est, notarum
eiusmodi tractationem, si calculos subeant prolixos termini, quos
adficiunt, et artificia, quibus ad servandam in illis concinnitatem opus
est, singulari disquisitione complecti.

1. Formulae propositae integratio initio quidem facillime perfici posse
videtur. Potest enim quantitas data $(1+n\cos\varphi)^\nu$ secundum
theorema binomiale in seriem per successivas ipsius $\cos\varphi$
potestates integras procedentem converti. Quo facto singuli eius termini,
in $d\varphi$ ducti, et ad integrationem revocati, quaesitum exhibebunt.
Recte hoc quidem, verum hac via incedendo ad expressiones complicatissimas
deducimur; quaelibet enim quaesitae integralis particula serie continebitur
haud parum prolixa. Praestat igitur aliam sequi rationem, in eo fundatam,
quod ex Trigonometriae analyticae praeceptis quaelibet ipsius $\cos\varphi$
potestas exponentis integri positivi possit in seriem valde concinnam, per
terminos qui solummodo cosinus multiplorum ipsius $\varphi$ continent,
progredientem evolvi. Ita igitur absolvetur negotium propositum, ut primo
quidem $(1+n\cos\varphi)^\nu$ per seriem, secundum potestates ipsius
$\cos\varphi$ procedentem, exhibeatur, mox vero, ex lege generali, antea
stabilita, quilibet eius terminus in seriem novam, per cosinus multiplorum
progredientem mutetur, atque colligantur demum ex omnibus hisce seriebus
membra, qui eiusdem anguli cosinum contineant. Series nova, eaque rite
ordinata, quae, peractis hisce laboribus exsurgit, in $d\varphi$ ducta
statim ad integrationem poterit perduci.

2. Ante omnia igitur lex, qua potestas aliqua cosinus dati per angulorum
multiplorum cosinus exhibetur, scrutanda erit. Huic disquisitioni
fundamento est formula trigonometrica elementaris\HistoricalNote{*)}{*)}{\textsc{Kaestner}, Astron. Abh., Erste Sammlung, I, 10.}, esse
